\documentclass[11pt]{article}
\usepackage[margin=2.5cm]{geometry}
\usepackage{amsmath,amssymb,amsthm}
\usepackage{microtype}

\newtheorem{lemma}{Lemma}
\newtheorem{proposition}{Proposition}
\theoremstyle{remark}
\newtheorem*{remark}{Remark}

\newcommand{\norm}[1]{\left\lVert #1 \right\rVert}
\newcommand{\Id}{I}

\title{Peaceman--Rachford against LOD for a diffusion step\\on a domain that changes every step}
\author{Development note (not for the paper)}
\date{}

\begin{document}
\maketitle

\section{Setting}

The viscous step of FLIP solves, on the current set of fluid faces,
\begin{equation}
  \partial_t u = \nu \nabla^2 u ,
\end{equation}
discretized in space as $u' = (A + B)\,u$, where $A$ collects the $x$ second differences and $B$ the $y$ ones. Let $k = \Delta t$ and
\begin{equation}
  r = \frac{\nu k}{h^2}
\end{equation}
be the diffusion number on a uniform grid of spacing $h$.

\paragraph{Unknowns.} The vector $u$ holds every face of the grid. A face that is not an unknown at step $n$ (solid, air, pinned) keeps its value: its row of $A_n$ and $B_n$ is zero, and no unknown couples to it. So $A_n$ and $B_n$ act on one fixed space, and only their entries change with $n$.

\paragraph{Properties used.} For each $n$:
\begin{enumerate}
  \item[(P1)] On a staircase domain with shared interface viscosities, $A_n$ and $B_n$ are symmetric and negative semidefinite: $A_n = A_n^{\mathsf T}$, $x^{\mathsf T} A_n x \le 0$.
  \item[(P2)] Their spectra lie in $[-4\nu/h^2,\,0]$, the bound of the one-dimensional second difference.
  \item[(P3)] In general $A_n B_n \ne B_n A_n$: the staircase breaks the tensor structure.
  \item[(P4)] On a fixed domain $A_n = A$, $B_n = B$; on FLIP's moving free surface they change at every step.
\end{enumerate}

\section{The two schemes}

\paragraph{Peaceman--Rachford (PR).} Two half steps, each implicit in one direction and explicit in the other:
\begin{align}
  \left(\Id - \tfrac{k}{2}A_n\right) u^{*} &= \left(\Id + \tfrac{k}{2}B_n\right) u^{n}, \\
  \left(\Id - \tfrac{k}{2}B_n\right) u^{n+1} &= \left(\Id + \tfrac{k}{2}A_n\right) u^{*},
\end{align}
so $u^{n+1} = G_n u^n$ with
\begin{equation}
  G_n = \left(\Id - \tfrac{k}{2}B_n\right)^{-1}\left(\Id + \tfrac{k}{2}A_n\right)\left(\Id - \tfrac{k}{2}A_n\right)^{-1}\left(\Id + \tfrac{k}{2}B_n\right).
  \label{eq:G}
\end{equation}

\paragraph{Locally one-dimensional (LOD).} Each sweep advances the full step in one direction, with no explicit operator:
\begin{equation}
  u^{n+1} = L_n u^n, \qquad L_n = \left(\Id - kB_n\right)^{-1}\left(\Id - kA_n\right)^{-1}.
  \label{eq:L}
\end{equation}

\section{Two elementary lemmas}

\begin{lemma}[Cayley factor]\label{lem:cayley}
Let $M$ be symmetric negative semidefinite and $k > 0$. Then
\begin{equation}
  C = \left(\Id + \tfrac{k}{2}M\right)\left(\Id - \tfrac{k}{2}M\right)^{-1}
\end{equation}
satisfies $\norm{C}_2 \le 1$.
\end{lemma}
\begin{proof}
$C$ is a function of the symmetric $M$, so it is symmetric with eigenvalues $(1 + k\mu/2)/(1 - k\mu/2)$, $\mu \le 0$ an eigenvalue of $M$. Each lies in $(-1, 1]$.
\end{proof}

\begin{lemma}[Implicit factor]\label{lem:implicit}
Let $M$ be symmetric negative semidefinite. Then $\norm{(\Id - kM)^{-1}}_2 \le 1$ for every $k \ge 0$.
\end{lemma}
\begin{proof}
Its eigenvalues are $1/(1 - k\mu) \in (0, 1]$ for $\mu \le 0$.
\end{proof}

\begin{remark}
The explicit factor has no such bound. $\Id + \tfrac{k}{2}M$ has eigenvalues $1 + k\mu/2$ with $\mu \in [-4\nu/h^2, 0]$ by (P2), so
\begin{equation}
  \norm{\Id + \tfrac{k}{2}M}_2 = \max\!\left(1,\; \left|1 - 2r\right|\right),
\end{equation}
which exceeds $1$ as soon as $r > 1$.
\end{remark}

\section{Fixed domain: PR is stable, without commutation}

\begin{proposition}\label{prop:fixed}
Let $A_n = A$ and $B_n = B$ for all $n$, both satisfying (P1)--(P2), and set $S = \Id - \tfrac{k}{2}B$. Then
\begin{equation}
  \norm{G^n}_2 \le \kappa(S) \le 1 + 2r \qquad \text{for all } n,
\end{equation}
whether or not $A$ and $B$ commute.
\end{proposition}
\begin{proof}
$\Id + \tfrac{k}{2}B$ and $S^{-1}$ are both functions of $B$, so they commute. From \eqref{eq:G},
\begin{equation}
  S\,G\,S^{-1}
  = \left(\Id + \tfrac{k}{2}A\right)\left(\Id - \tfrac{k}{2}A\right)^{-1} \left(\Id + \tfrac{k}{2}B\right) S^{-1}
  = C_A\, C_B ,
\end{equation}
with $C_A$, $C_B$ the Cayley factors of Lemma~\ref{lem:cayley}. Hence $\norm{SGS^{-1}}_2 \le 1$, and the similarity telescopes:
\begin{equation}
  G^n = S^{-1} \left(C_A C_B\right)^n S
  \quad\Longrightarrow\quad
  \norm{G^n}_2 \le \norm{S^{-1}}_2 \norm{S}_2 = \kappa(S).
\end{equation}
By (P2) the eigenvalues of $S$ lie in $[1,\, 1 + 2r]$, so $\kappa(S) \le 1 + 2r$.
\end{proof}

\begin{remark}
PR is a contraction in the norm $\norm{x}_S = \norm{Sx}_2$. In the plain norm a transient growth up to $1 + 2r$ is possible, but it never accumulates. Commutation of $A$ and $B$ is not needed; it would only make $G$ itself a contraction in $\norm{\cdot}_2$.
\end{remark}

\section{Moving domain: the similarity no longer cancels}

Now $A_n$, $B_n$ change at every step (P4). Let $S_n = \Id - \tfrac{k}{2}B_n$ and $C_n = C_{A_n} C_{B_n}$. By the same algebra, $G_n = S_n^{-1} C_n S_n$, and the product of $N$ steps is
\begin{multline}
  G_{N-1}\cdots G_1 G_0 \\
  = S_{N-1}^{-1}\, C_{N-1}\, \bigl(S_{N-1} S_{N-2}^{-1}\bigr)\, C_{N-2} \cdots \bigl(S_1 S_0^{-1}\bigr)\, C_0\, S_0 .
  \label{eq:product}
\end{multline}
The interior factors $S_{j} S_{j-1}^{-1}$ are the identity only when $B_j = B_{j-1}$. Since $\norm{C_j}_2 \le 1$, and since $\norm{S_j}_2 \norm{S_{j-1}^{-1}}_2 \le 1 + 2r$ for every $j$,
\begin{equation}
  \norm{G_{N-1}\cdots G_0}_2 \le (1 + 2r)^{N}.
  \label{eq:bound}
\end{equation}

\begin{remark}[What this does and does not prove]
\eqref{eq:bound} is an upper bound: it shows that the argument guaranteeing PR's stability on a fixed domain fails on a moving one, and that each change of geometry can amplify by up to $1 + 2r$. It does not prove that growth occurs; that depends on how the domain changes. For $r \lesssim 1$ the bound is mild, and for $r \gg 1$, FLIP's viscous regime, it allows exponential growth. The mechanism is the explicit factor $\Id + \tfrac{k}{2}B_n$: its amplification, of norm up to $|1 - 2r|$, is undone only by the implicit factor $(\Id - \tfrac{k}{2}B_n)^{-1}$ with the \emph{same} $B_n$, which the next step no longer applies.
\end{remark}

\section{LOD is a contraction for every sequence of domains}

\begin{proposition}\label{prop:lod}
Under (P1), for any $k \ge 0$ and any sequence $(A_n, B_n)$,
\begin{equation}
  \norm{L_{N-1}\cdots L_0}_2 \le 1 .
\end{equation}
\end{proposition}
\begin{proof}
Each factor of \eqref{eq:L} has $\norm{\cdot}_2 \le 1$ by Lemma~\ref{lem:implicit}, and the norm is submultiplicative. No similarity transform is involved, so nothing depends on consecutive steps sharing an operator.
\end{proof}

\section{Without symmetry: the maximum norm}

FLIP's assembled viscous matrices are not exactly symmetric: the free-surface gating of a tangential coupling is decided from each row's own face, so a pair of rows can disagree. (P1) then fails and the $2$-norm arguments above do not apply. A weaker structure survives: for each row $i$ of $M = A_n$ or $B_n$,
\begin{equation}
  m_{ii} = -d_i \le 0, \qquad m_{ij} \ge 0 \ (j \ne i), \qquad \sum_{j \ne i} m_{ij} \le d_i ,
  \label{eq:mmatrix}
\end{equation}
that is, $-M$ is a weakly diagonally dominant $M$-matrix. A solid neighbour keeps its diagonal term and drops the coupling, which only strengthens the dominance.

\begin{lemma}\label{lem:maxnorm}
Under \eqref{eq:mmatrix}, $\norm{(\Id - kM)^{-1}}_\infty \le 1$ for every $k \ge 0$.
\end{lemma}
\begin{proof}
$T = \Id - kM$ has diagonal $1 + kd_i > 0$, non-positive off-diagonal entries, and row sums $1 + k(d_i - \sum_{j\ne i} m_{ij}) \ge 1$. It is a strictly diagonally dominant $M$-matrix, so $T^{-1} \ge 0$ entrywise. With $\mathbf{1}$ the vector of ones, let $y = T^{-1}\mathbf{1} \ge 0$ (the row sums of $T^{-1}$) and $y_p = \max_i y_i$. Row $p$ of $Ty = \mathbf{1}$ reads
\begin{equation}
  1 = (1 + kd_p)\, y_p - k\sum_{j\ne p} m_{pj}\, y_j \;\ge\; (1 + kd_p)\, y_p - k d_p\, y_p = y_p .
\end{equation}
So every row sum of the non-negative $T^{-1}$ is at most $1$, which is $\norm{T^{-1}}_\infty \le 1$.
\end{proof}

\begin{proposition}
Under \eqref{eq:mmatrix}, LOD satisfies $\norm{L_{N-1}\cdots L_0}_\infty \le 1$ for every $k$ and every sequence of domains. The PR explicit factor, at an interior row with $d_i = 2\nu/h^2$ and $\sum_{j\ne i} m_{ij} = d_i$, has
\begin{equation}
  \norm{\Id + \tfrac{k}{2}M}_\infty \;\ge\; \left|1 - r\right| + r ,
\end{equation}
which is $2r - 1 > 1$ for $r > 1$.
\end{proposition}
\begin{proof}
The LOD bound follows from Lemma~\ref{lem:maxnorm} and submultiplicativity. For the explicit factor, row $i$ of $\Id + \tfrac{k}{2}M$ has diagonal $1 - \tfrac{k}{2}d_i = 1 - r$ and off-diagonal sum $\tfrac{k}{2}d_i = r$, all off-diagonal entries non-negative.
\end{proof}

\begin{remark}
Without symmetry there is also no norm in which each Cayley factor is a contraction for all $n$ at once, so even on a fixed domain the PR argument of Proposition~\ref{prop:fixed} is not available; for FLIP, which has both the moving domain and the asymmetry, nothing bounds PR, while LOD stays a contraction in $\norm{\cdot}_\infty$.
\end{remark}

\section{Literature}

\paragraph{Hundsdorfer and Verwer (1989)}, \emph{Math.\ Comp.} 53:81--101, Theorem~3.2. For dissipative, possibly non-commuting and time-dependent operators, they write PR with the explicit $F_2$ of step $n$ evaluated at $t_n$, the same time level as the implicit $F_2$ that ended step $n-1$. Their transformation $\varepsilon_n \mapsto (\Id - \tfrac{\tau}{2}A_{2,n})\varepsilon_n$ is our $S_n$, and because both factors use the same $A_{2,n}$ the similarity telescopes exactly. This is Proposition~\ref{prop:fixed} for time-dependent operators. The only loose end is the first step, whose transformed error $(\Id - \tfrac{\tau}{2}A_{2,0})\varepsilon_0$ ``may become quite large and grow with spatial refinement''. Their remedy is one backward Euler LOD step first. They also note that PR damps high-frequency errors only slowly, while LOD damps them strongly.

\paragraph{The link to FLIP.} The free surface is re-marked between steps, so the operator of the explicit half step of step $n$ is not the one of the implicit half step that ended step $n-1$. That is the factor $S_n S_{n-1}^{-1}$ of \eqref{eq:product}. In Hundsdorfer and Verwer's terms every FLIP step is a new first step, with a nonsmooth error, and their remedy applied every step is LOD throughout.

\paragraph{Related.} Schatzman (1999), \emph{J.\ Funct.\ Anal.} 162:219--255, shows that for non-commuting unbounded self-adjoint operators the PR product is not stable in general, the operator-level counterpart of $\kappa(S) \sim 1 + 2r$ growing without bound as $h \to 0$ (not read beyond the abstract). Li and Zhao (2017), \emph{Appl.\ Math.\ Comput.} 299:28--44, report that PR loses unconditional stability at interfaces where Douglas ADI keeps it (not read).

\section{Summary and predictions}

\begin{center}
\begin{tabular}{lcc}
  & fixed domain & domain changing every step \\[2pt]
  \hline
  PR, symmetric operators & bounded by $1 + 2r$ & bound $(1+2r)^N$, growth possible for $r > 1$ \\
  PR, non-symmetric (FLIP) & no bound available & no bound available \\
  LOD, symmetric & $\norm{\cdot}_2 \le 1$ & $\norm{\cdot}_2 \le 1$ \\
  LOD, $M$-matrix (FLIP) & $\norm{\cdot}_\infty \le 1$ & $\norm{\cdot}_\infty \le 1$ \\
\end{tabular}
\end{center}

\noindent Predictions to test on a 2D staircase diffusion problem, sweeping $r$ from $0.1$ to $100$:
\begin{enumerate}
  \item PR on a fixed staircase mask: $\norm{u^n}$ bounded, possibly with transient growth up to $1 + 2r$.
  \item PR on a mask that changes every step: growth possible once $r > 1$.
  \item LOD: $\norm{u^n}_2$ and $\norm{u^n}_\infty$ non-increasing in every case.
\end{enumerate}

\noindent The price of LOD is accuracy: its splitting error is first order in $k$, and by Weare (1979) any splitting carries the boundary error into the interior at a rate growing with the time step. For FLIP, whose other errors are first order already, stability is the property worth buying.

\end{document}
